\heading{实验物理中的统计方法-模拟题3-期中}
\noteinfo{题目和答案由AI生成。知识点范围按现有往年题整理，叙述风格尽量向往年题靠拢；本套难度较前两套期中模拟题再提高一些。}

\begin{question}
从 $\{1,2,\cdots,n\}$ 中随机取一个数 $X$，再从 $\{1,2,\cdots,X\}$ 中随机取一个数 $Y$。
\begin{enumerate}[label=(\alph*),leftmargin=2.8em,itemsep=0.2em]
\item 求 $P(Y=m)$（其中 $1\le m\le n$）；
\item 求 $E[Y]$。
\end{enumerate}
\begin{solution}
对 $1\le m\le n$，有
\[
P(Y=m)=\sum_{x=m}^{n}P(X=x)P(Y=m\mid X=x)
=\frac1n\sum_{x=m}^{n}\frac1x.
\]
因此
\ansmath{P(Y=m)=\frac1n\sum_{x=m}^{n}\frac1x=\frac{H_n-H_{m-1}}{n}}
其中 $H_k=\sum_{j=1}^k\frac1j$ 为调和数。
又由条件期望公式，
\[
E[Y\mid X=x]=\frac{x+1}{2},
\qquad
E[Y]=E\bigl(E[Y\mid X]\bigr)=\frac{E[X]+1}{2}.
\]
而
\[
E[X]=\frac{n+1}{2}.
\]
\end{solution}
\end{question}

\begin{question}
设 $X,Y$ 相互独立，且都服从 $U(0,1)$。令
\[
U=\min(X,Y),\qquad V=\max(X,Y).
\]
\begin{enumerate}[label=(\alph*),leftmargin=2.8em,itemsep=0.2em]
\item 求 $(U,V)$ 的联合概率密度；
\item 求 $P(V-U<t)$，其中 $0<t<1$。
\end{enumerate}
\begin{solution}
在区域 $0<u<v<1$ 上，$(u,v)$ 对应 $(x,y)$ 的两个排列，因此联合密度为
\[
f_{U,V}(u,v)=2,
\qquad 0<u<v<1.
\]
故
\ansmath{f_{U,V}(u,v)=\begin{cases}2,&0<u<v<1,\\0,&\text{其他} .\end{cases}}
再求
\[
P(V-U<t)=P(|X-Y|<t).
\]
在单位正方形中，满足 $|x-y|\ge t$ 的区域是两块边长为 $1-t$ 的直角三角形，总面积为 $(1-t)^2$。因此
\[
P(|X-Y|<t)=1-(1-t)^2=2t-t^2.
\]
\end{solution}
\end{question}

\begin{question}
某电路由三个相互独立、同分布的元件 $A,B,C$ 组成。只有当元件 $A$ 正常，且元件 $B,C$ 至少有一个正常时，整个电路才正常。设单个元件正常工作时间的累积分布函数为 $F(t)$。求整个电路正常工作时间 $T$ 的累积分布函数。
\begin{solution}
电路在时刻 $t$ 仍正常工作的条件是：
\[
A>t,\qquad B>t\ \text{或}\ C>t.
\]
因此
\[
P(T>t)=P(A>t)P(B>t\ \text{或}\ C>t).
\]
又
\[
P(A>t)=1-F(t),
\]
\[
P(B>t\ \text{或}\ C>t)=1-P(B\le t,C\le t)=1-F(t)^2.
\]
故
\[
P(T>t)=[1-F(t)][1-F(t)^2].
\]
\end{solution}
\end{question}

\begin{question}
随机变量 $(X,Y)$ 的联合概率密度为
\[
f(x,y)=\begin{cases}
 c(x+2y),&0<y<x<1,\\
 0,&\text{其他}.
\end{cases}
\]
\begin{enumerate}[label=(\alph*),leftmargin=2.8em,itemsep=0.2em]
\item 求常数 $c$；
\item 求边缘密度 $f_X(x)$；
\item 求 $P\!\left(Y>\dfrac{X}{2}\right)$；
\item 求 $E[Y\mid X=x]$。
\end{enumerate}
\begin{solution}
由归一化条件，
\[
1=c\int_0^1\int_0^x (x+2y)\,dy\,dx
=c\int_0^1 2x^2\,dx
=\frac{2c}{3}.
\]
故
\ansmath{c=\frac32}
边缘密度为
\[
f_X(x)=\int_0^x \frac32(x+2y)\,dy
=\frac32\cdot 2x^2=3x^2,
\qquad 0<x<1.
\]
所以
再算概率：
\[
P\!\left(Y>\frac X2\right)=\int_0^1\int_{x/2}^{x}\frac32(x+2y)\,dy\,dx.
\]
计算得
\[
\int_{x/2}^{x}(x+2y)\,dy=\frac{5x^2}{4},
\]
故
\[
P\!\left(Y>\frac X2\right)=\int_0^1 \frac32\cdot \frac{5x^2}{4}\,dx=\frac58.
\]
条件密度为
\[
f_{Y\mid X}(y\mid x)=\frac{\frac32(x+2y)}{3x^2}=\frac{x+2y}{2x^2},
\qquad 0<y<x.
\]
于是
\[
E[Y\mid X=x]=\int_0^x y\,\frac{x+2y}{2x^2}\,dy=\frac{7x}{12}.
\]
\end{solution}
\end{question}

\begin{question}
设 $\{N(t),t\ge 0\}$ 是参数为 $\lambda$ 的泊松过程，$T_1<T_2<T_3<\cdots$ 为各次到达时刻。求条件密度 $f_{T_2\mid T_3=t}(u)$。
\begin{solution}
条件于 $T_3=t$ 后，$T_1,T_2$ 就是在区间 $(0,t)$ 上两个独立均匀点的顺序统计量。因此
\[
\frac{T_2}{t}
\]
服从两个 $U(0,1)$ 样本中较大者的分布，其密度为
\[
f(w)=2w,
\qquad 0<w<1.
\]
令 $w=u/t$，则
\[
f_{T_2\mid T_3=t}(u)=2\frac{u}{t}\cdot \frac1t
=\frac{2u}{t^2},
\qquad 0<u<t.
\]
\end{solution}
\end{question}

\begin{question}
用宇宙射线标定探测器。把两块完全相同、相互独立的待测闪烁体串联放在上下两块触发闪烁体之间。现记录到“上下两块都击中”的有效穿过共 $3000$ 次，其中两块待测闪烁体都击中的次数为 $2523$ 次。
\begin{enumerate}[label=(\alph*),leftmargin=2.8em,itemsep=0.2em]
\item 设单块待测闪烁体效率为 $\varepsilon$，求 $\varepsilon$ 的估计值；
\item 求该估计值的统计不确定度；
\item 若把这两块闪烁体改成“只要至少一块击中就算通过”，求系统效率的估计值。
\end{enumerate}
\begin{solution}
设两块都击中的概率为 $p=\varepsilon^2$。由数据得
\[
\hat p=\frac{2523}{3000}=0.841.
\]
故
\[
\hat\varepsilon=\sqrt{\hat p}=\sqrt{0.841}\approx 0.9171.
\]
按二项分布近似，
\[
\sigma_{\hat p}=\sqrt{\frac{\hat p(1-\hat p)}{N}}
=\sqrt{\frac{0.841\times 0.159}{3000}}
\approx 0.00668.
\]
又因 $\varepsilon=\sqrt p$，故
\[
\sigma_{\hat\varepsilon}\approx \left|\frac{d\varepsilon}{dp}\right|\sigma_{\hat p}
=\frac{1}{2\hat\varepsilon}\sigma_{\hat p}
\approx \frac{0.00668}{2\times 0.9171}
\approx 0.00364.
\]
因此
\ansmath{\hat\varepsilon=0.917\pm 0.004}
若改成“至少一块击中”，则系统效率为
\[
\varepsilon_{\rm sys}=1-(1-\varepsilon)^2.
\]
代入估计值得
\[
\hat\varepsilon_{\rm sys}=1-(1-0.9171)^2\approx 0.9931.
\]
\end{solution}
\end{question}

\begin{question}
设探测效率 $\varepsilon$ 的先验分布为
\[
\pi(\varepsilon)\propto \varepsilon(1-\varepsilon)^2,
\qquad 0<\varepsilon<1.
\]
现做了 $10$ 次标准刻度，其中有 $8$ 次探测到信号。
\begin{enumerate}[label=(\alph*),leftmargin=2.8em,itemsep=0.2em]
\item 求后验分布；
\item 求后验均值；
\item 求接下来两次刻度都探测到信号的后验预测概率。
\end{enumerate}
\begin{solution}
先验可写成
\[
\varepsilon\sim \mathrm{Beta}(2,3).
\]
观测到 $8$ 次成功、$2$ 次失败后，后验分布为
\[
\varepsilon\mid D\sim \mathrm{Beta}(10,5).
\]
后验均值为
\[
E[\varepsilon\mid D]=\frac{10}{15}=\frac23.
\]
后面两次都成功的后验预测概率为
\[
E[\varepsilon^2\mid D]
=\frac{10\cdot 11}{15\cdot 16}
=\frac{11}{24}.
\]
\end{solution}
\end{question}

\begin{question}
某稀有事件搜索中，本底计数服从参数为 $3$ 的泊松分布。现观测到 $7$ 个候选事例。若只考虑“本底涨落到不小于观测值”的显著性水平，求对应的 $p$ 值。
\begin{solution}
所求为
\[
p=P(N\ge 7\mid \lambda=3)=1-P(N\le 6\mid \lambda=3).
\]
即
\[
p=1-e^{-3}\sum_{k=0}^{6}\frac{3^k}{k!}.
\]
数值上
\[
p\approx 0.0335.
\]
\end{solution}
\end{question}

